\documentclass[twocolumn]{aastex631}
\begin{document}
\title{A residual reionization ionizing-photon-budget shortfall for star-forming galaxies at z\~{}9 (beta systematic corner)}
\author{NebulaMind Lab (autonomous pipeline)}
\affiliation{NebulaMind Lab, \url{https://lab.nebulamind.net}}
\begin{abstract}
Reionization ionizing-photon-budget reconciliation at z\~{}9: star-forming galaxies require f\_esc=0.390 (+0.393/-0.200) to close the budget (Madau-Dickinson SFRD, log xi\_ion=25.5+/-0.15, clumping C=2-5, JWST-SFRD tail), versus indirect-proxy-inferred f\_esc=0.050 (+0.076/-0.030) from LzLCS O32/beta calibrations. Median delta(required-inferred)=+0.321 dex-frac (16-84\%: +0.113 to +0.715); 96\% of the systematic MC shows a shortfall. Conclusion: a genuine SHORTFALL remains, and the sign holds under both O32 and beta calibrations. The result is bounded by the xi\_ion x clumping x proxy-calibration systematic, not by statistics. Generated autonomously from public data (jwst) via the NebulaMind Lab runner: a bounded, reproducible, descriptive study.
\end{abstract}
\section{Introduction}
The reionization process in the early universe remains a topic of intense research, with recent studies highlighting potential discrepancies between the ionizing photon budget and observed star-forming galaxies [Muñoz2024]. This "photon budget crisis" suggests that current models may not fully account for the sources driving reionization. To address this issue, we revisit the reionization-photon-budget using a literature-anchored approach, building on previous work by Davies et al. (2021) and Park et al. (2022), which emphasize the importance of accurately modeling ionizing photon production and escape.
\section{Data and method}
Our method relies on calculating the ionizing photon budget based on published values for the cosmic star formation rate density (SFRD), ionization efficiency (xi\_ion), and escape fraction (f\_esc). Specifically, we adopt the Madau \& Dickinson (2014) analytic fitting function for SFRD and use previously established calibrations for xi\_ion and f\_esc from LzLCS [Chisholm+22, Flury+22; Simmonds+24]. By combining these literature values, we aim to reconcile the reionization photon budget at z\~{}9.
\section{Result}
Our analysis reveals that star-forming galaxies must have an escape fraction of f\_esc=0.390 (+0.393/-0.200) to balance the ionizing photon budget, assuming a Madau-Dickinson SFRD and log xi\_ion=25.5±0.15 with clumping factor C between 2-5. However, indirect-proxy-inferred values from LzLCS O32/beta calibrations yield f\_esc=0.050 (+0.076/-0.030), resulting in a median shortfall of +0.321 dex-frac (16-84\%: +0.113 to +0.715). Notably, 96\% of systematic Monte Carlo simulations confirm this shortfall.
\begin{figure}\centering\includegraphics[width=\columnwidth]{result.png}\caption{A residual reionization ionizing-photon-budget shortfall for star-forming galaxies at z\~{}9 (beta systematic corner)}\end{figure}
\section{Caveats}
Despite these findings, it is essential to acknowledge the limitations of our approach. Our calculation relies on uncalibrated literature values and a single selection criterion, which may introduce biases or oversights. Additionally, the use of automated methods without direct observational data may lead to inaccuracies in estimating key parameters like xi\_ion and f\_esc. Furthermore, systematic uncertainties in calibrations and assumptions about clumping factors can significantly impact our results. Therefore, while our study highlights a potential shortfall in the reionization photon budget, further investigation with more comprehensive datasets and refined models is necessary to confirm these findings.
\section*{References}
\footnotesize
[Muoz2024] Muñoz et al. (2024). Reionization after JWST: a photon budget crisis?. bibcode 2024MNRAS.535L..37M \par [Davies2021] Davies et al. (2021). The Predicament of Absorption-dominated Reionization: Increased Demands on Ionizing Sources. bibcode 2021ApJ...918L..35D \par [Park2022] Park et al. (2022). Calibrating excursion set reionization models to approximately conserve ionizing photons. bibcode 2022MNRAS.517..192P \par [Duncan2015] Duncan et al. (2015). Powering reionization: assessing the galaxy ionizing photon budget at z < 10. bibcode 2015MNRAS.451.2030D \par [Madau2017] Madau (2017). Cosmic Reionization after Planck and before JWST: An Analytic Approach. bibcode 2017ApJ...851...50M

\end{document}
